\chapter*{Введение}
\addcontentsline{toc}{chapter}{Введение}\markboth{Введение}{Введение}
\bookpage{5}

Существует обширный круг научно-технических проблем, при решении которых
необходимо учитывать и преобразовывать геометрическую информацию.
В простейших случаях (например, при задании информации о машиностроительном
чертеже для хранения ее в памяти ЭВМ) вполне можно ограничиться тем или иным
способом кодирования элементов чертежа или в крайнем случае поточечным их
сканированием. Такие задачи в большинстве случаев носят чисто технический
характер и их решение не сопряжено с преодолением принципиальных затруднений.
Вместе с тем во многих случаях простого кодирования геометрической
информации недостаточно, а необходимо описание ее в той форме, какая принята
в аналитической геометрии, т.\,е. с помощью уравнений или неравенств. Речь
идет о том, чтобы так же, как, например, для прямой на плоскости или шара
в пространстве, можно было бы и для таких <<сложных>> геометрических
объектов, как усеченный конус с цилиндрическим отверстием, ломаная линия,
зубчатое колесо и других применять уравнения $\omega(x)=0$, $x\in\mathbb R^n$
(или неравенства $\omega(x)\ge0$). Здесь $\omega(x)$ "--- функция, которая
имеет вид единого аналитического выражения и обладает, возможно, еще
и некоторыми дополнительными свойствами.

Естественно, что при такой общей постановке возникает ряд вопросов,
требующих строгой формализации (и такая формализация выполнена в данной
работе). К их числу относятся следующие: что такое сложный геометрический
объект и единое аналитическое выражение, при каких условиях поставленная
задача разрешима, имеются ли эффективные алгоритмы решения этой задачи
и~т.\,д.?

В историческом плане сформулированная выше проблема восходит еще к Декарту
и известна в литературе как обратная задача аналитической геометрии: задан
геометрический объект, требуется написать его уравнение. (Прямая задача:
дано уравнение $\omega(x)=0$, $x\in\mathbb R^n$, найти соответствующий ему
в~$\mathbb R^n$ геометрический объект $L$, т.\,е. совокупность всех точек
$\mathbb R^n$, которые удовлетворяют уравнению $\omega(x)=0$.) Во времена
Декарта запас изученных линий и поверхностей (чертежей) мало чем отличался от
того, каким располагали математики древности. Метод координат
\bookpage{6}%
позволил беспредельно увеличивать число изученных чертежей, так как каждое
новое уравнение давало и новый чертеж. Это привело к тому, что в дальнейшем
основное внимание было уделено изучению чертежей, описываемых заданными
уравнениями. Обратная же задача рассматривалась лишь для простейших
геометрических объектов: прямой, конических сечений, эллипсоидов,
параболоидов и~др.

В аналитической геометрии обычно ограничиваются системой
$H_0=\{x_1+x_2,\ x_1x_2,\ a=\mathrm{const}\ \forall a\in\mathbb R\}$
операций, применяемых для написания уравнений, которая приводит к множеству
$\mathfrak M(H_0)$ целых рациональных функций (полиномов). Ему соответствует
множество $\mathfrak R(H_0)$ чертежей, описываемых уравнениями $\omega(x)=0$,
$\omega(x)\in\mathfrak M(H_0)$ и называемых алгебраическими. Однако множество
алгебраических чертежей все же весьма <<бедно>>. Оно не содержит, например,
таких чертежей, как прямоугольник, двутавр, поверхность искусственного
спутника Земли и~т.\,п. Правда, существует принципиальная возможность
аппроксимации сложных чертежей алгебраическими (это следует из известной
теоремы Вейерштрасса о приближении непрерывных функций полиномами~[7]), но
сама по себе реализация ее весьма сложная проблема. То же можно сказать
и относительно известной теоремы о возможности сколь угодно хорошей
аппроксимации замкнутых кривых на плоскости лемнискатами (т.\,е. также
алгебраическими кривыми)~[150].

Существенным расширением множества $\mathfrak R(H_0)$ алгебраических
чертежей является множество $\mathfrak R_{1/2}(H_0)$ полуалгебраических
чертежей~[5]. Оно включает в себя всякие чертежи, которые также могут быть
описаны с помощью упомянутой системы операций $H_0$, но уже не одним
уравнением, а системами систем уравнений и неравенств вида $f_{ik}=0$,
$\varphi_{jk}\ge0$, где $f_{ik},\varphi_{jk}\in\mathfrak M(H_0)$, $i$, $j$,
$k$ "--- принадлежат конечным множествам индексов (точка принадлежит
чертежу, если хотя бы одна из систем удовлетворяется). Из определения
множества $\mathfrak R_{1/2}(H_0)$ следует, что
$\mathfrak R(H_0)\subset\mathfrak R_{1/2}(H_0)$. Нетрудно установить также
(см. \S\,9 гл.~2), что существуют полуалгебраические, но не алгебраические
чертежи, т.\,е. $\mathfrak R(H_0)\ne\mathfrak R_{1/2}(H_0)$.

Пусть $H$ есть некоторое расширение системы операций $H_0$. Тогда, подобно
тому как системе $H_0$ были поставлены в соответствие множества чертежей
$\mathfrak R(H_0)$ и $\mathfrak R_{1/2}(H_0)$, системе $H$ будут
соответствовать множества $\mathfrak R(H)$ и $\mathfrak R_{1/2}(H)$
чертежей, которые назовем $H$-реализуемыми и полу-$H$-реализуемыми
соответственно. Правомерен вопрос: нельзя ли так выбрать систему $H$, чтобы
получить равенство $\mathfrak R(H)=\mathfrak R_{1/2}(H)$? Располагая такой
системой, можно было бы задавать сложные геометрические объекты не
упомянуты\-%
\bookpage{7}%
ми выше системами систем уравнений и неравенств, а одним уравнением,
написанным в символах системы $H$. В гл.~2 будет показано, что такие
системы, названные алгоритмически полными, существуют и могут быть
построены с помощью $R$-функций. Впервые понятие $R$-функции как раз и было
введено в связи с возникшей необходимостью составлять уравнения сложных
чертежей (в процессе поиска приближенных решений некоторых пространственных
контактных задач теории упругости).

Со временем понятие $R$-функции было обобщено и возникло общее понятие
$R$-отображения (см. \S\,2 гл.~2). В основе этого понятия лежит выделение из
систем алгебраических операций (т.\,е. отображений вида
$f\colon\mathfrak X^n\to\mathfrak X$, где $\mathfrak X$ "--- некоторое
множество) таких, которые в некотором смысле соответствуют функциям
$k$-значной логики (т.\,е. алгебраическим операциям вида
$F\colon B_k^n\to B_k$, у которых множество $B_k$ состоит из $k$
элементов). В общем виде теория $R$-отображений еще мало разработана.
В основном изучались $R$-функции, которые строятся на множестве
$\mathfrak X=\mathbb R=(-\infty,\infty)$, и соответствуют функциям
двузначной (булевым функциям) и трехзначной логики. Именно эти функции
позволили решить обратную задачу аналитической геометрии, а впоследствии
найти и другие приложения на более высоких уровнях прикладной математики.
При этом оказалось, что существуют элементарные алгоритмически полные
системы и в целом само множество основных элементарных функций (в сочетании
с арифметическими операциями) алгоритмически полное. Это означает,
в частности, что всякий чертеж, составленный из <<кусков>> элементарных
линий или поверхностей, сам элементарен, т.\,е. может быть описан
уравнением $\omega(x)=0$, где $\omega(x)$ "--- элементарная (имеет вид
единого аналитического выражения) функция.

В первых работах по теории $R$-функций был определен и достаточно широкий
круг возможных ее приложений: оптимальное размещение геометрических объектов
(в частности, оптимальный раскрой [109, 142, 143]), распознавание образов
[157, 158], математическое программирование [109], конструктивная теория
функций [61, 111] и в особенности краевые задачи математической физики [111,
128, 134].

Одно из неожиданных направлений, в котором $R$-функции нашли применение,
"--- это теория устойчивости движения. Первоначально метод $R$-функций
применен в этом направлении при исследовании области устойчивости
неустановившегося движения~[77]. За истекшее время с использованием
$R$-функций получены новые результаты при численной реализации классических
критериев устойчивости [78, 79, 82], построении функций Ляпунова, синтезе
систем, асимптотически устойчивых по первому приближению~[79], выборе
параметров функций связей крупномасштабных систем с лишь устойчивыми
подсистемами~[80], в задаче об устойчивости аналитических движений~[76]
и других~[77]. Этим, судя по всему, не исчерпываются возможности применения
$R$-функций. В част\-%
\bookpage{8}%
ности, представляют значительный интерес разработка теории
$R$-отображений в пространстве параметров динамических систем~[81],
построение и исследование свойств $R$-отображений, управляемых специальными
системами сравнения [76--82], а также применение $R$-функций в задачах
распознавания образов [157, 158], химической технологии~[46], медицинской
диагностики~[19] и других направлениях [7, 36].

Математическому программированию и размещению геометрических объектов
посвящены \S\S\,26, 27 гл.~2. В основном они написаны на уровне
монографии~[109], изданной в 1967~г., более глубоко эти вопросы освещены
в работах Ю.\,Г.~Стояна и его учеников [167--172]. Кратко изложены также
вопросы, относящиеся к теории сплайнов и атомарных функций (\S\S\,15, 16
гл.~3). Более детально эти вопросы рассмотрены в монографии~[127]. Основные
результаты в развитии теории атомарных функций принадлежат В.\,А.~Рвачеву
[123--127], первые приложения "--- Т.\,И.~Шейко~[184].

В последние годы в индустриально развитых странах возросли темпы создания
новых математических теорий и разработки на их основе методов решения на ЭВМ
задач исследования, расчета и оптимизации физико-механических полей. Все
возрастающий интерес к проблеме физико-механических полей объясняется тем,
что ею охвачен широкий круг направлений, таких как электродинамика,
теплофизика, теория фильтрации, теория упругости и пластичности, магнитная
гидродинамика и другие, развитие которых имеет первостепенное значение для
научно-технического прогресса. Математическими моделями физико-механических
полей являются задачи для уравнений с частными производными при
определенных краевых и начальных условиях. Специфическая особенность полей
"--- их зависимость не только от характера физических законов, учитываемых
соответствующими уравнениями, но и от формы, взаимного расположения тел,
в которых поля возбуждаются, конфигурации площадок их контактного
взаимодействия и других геометрических и физических факторов.

Наличие в постановке краевых задач двух разнородных видов информации "---
аналитической и геометрической "--- серьезное препятствие при создании
методов и алгоритмов решения этих задач. Дело в том, что всякий метод
неизбежно должен учитывать оба упомянутые вида информации, а это требует
преобразования геометрической информации в аналитическую (при этом не всякое
кодирование геометрической информации годится). В таких классических
методах, как, например, разделения переменных и интегральных преобразований,
форма областей и участков их границ учитывается благодаря подходящему выбору
систем координат, в методе конформных отображений "--- при построении
отображающих функций, в вариационных методах "--- при построении
координатных функций, в сеточных методах "--- при составлении уравнений
для узлов, близких к границе, и~т.\,д. Метод конечных элементов
(МКЭ)~[31, 66] возник именно в связи со стремлением как
\bookpage{9}%
можно точнее учитывать геометрические формы. Появился также ряд методов
(Рейснера, штрафов и~др.), в которых краевые условия естественны и
формально им можно не удовлетворять [58, 88]. В этих методах геометрическая
информация учитывается путем интегрирования по областям соответствующей
формы (т.\,е. геометрическая информация учитывается интегрально).

Особенно актуальна разработка таких методов решения краевых задач, которые
имели бы универсальный характер и не требовали от исследователя (как
правило, не математика) знания тонких вопросов теории. Кроме того,
универсальность обусловливает возможность привлечения методов системного
программирования, что имеет существенное значение для автоматизации научных
исследований в области краевых задач. Конечно, для того или иного
конкретного класса задач существуют или могут быть созданы специальные
методы, превосходящие по эффективности любой универсальный метод, но их
использование требует, как правило, хорошей математической и специальной
подготовки и значительных затрат времени, что может оказаться неприемлемым
для быстрого решения все новых и новых задач, различных по специфике. Из
упомянутых выше методов универсальны разностные, вариационные и конечных
элементов. Можно назвать также методы Трефтца, локальных вариаций~[183],
штрафов~[112] и~др. Перечисленные методы принято называть прямыми; для них
характерно сведение краевых задач к решению алгебраических систем линейных
уравнений. Есть, однако, другая, более глубокая связь между этими методами,
которая делает различие между ними весьма условным: во всех прямых методах
приближенное решение краевой задачи отыскивается в виде
\begin{equation}\tag{1}\label{eq:intro1}
u_N=\sum_{k=1}^{N}C_k\varphi_k(x)+\varphi_0(x),
\end{equation}
где $\varphi_k(x)$ "--- некоторые заранее выбранные функции, а $C_k$ "---
числа, которые находятся в результате решения упоминавшихся выше систем
алгебраических уравнений. Заметим, что в некоторых методах (сеточных, МКЭ)
формула~(1) представления решения не всегда явно выписывается, однако это
не меняет сути. В простейших конечно-разностных методах, например, за основу
берется функция $y=h_1(t)$ (финитный сплайн Шенберга первого порядка).

Приближенное решение (например, для двумерной задачи) выбирается в виде
линейной комбинации произведений сдвигов сжатий функций~$h_1$:
\begin{equation}\tag{2}\label{eq:intro2}
u_N(x_1,x_2)=\sum_{i,j}C_{ij}\,h_1(\alpha x_1-\beta i)\,h_1(\gamma x_2-\delta j).
\end{equation}
Здесь $C_{ij}=u_{ij}$ "--- значение искомой функции в узле $(i,j)$, а
$\alpha$, $\beta$, $\gamma$ и $\delta$ выбираются такими, чтобы функция
$\varphi_{ij}=h_1(\alpha x_1-\beta i)\,h_1(\gamma x_2-\delta j)$ была равна
единице в $(i,j)$-м узле и
\bookpage{10}%
нулю "--- в остальных узлах сетки. Нетрудно представить, что графически
функция~(2) "--- это многогранная поверхность. При этом изменение
какой-либо из величин $C_{ij}$ приводит лишь к локальному ее изменению
в окрестности узла $(i,j)$. Если, учтя граничные условия, подставить
формулу~(2) в функционал, соответствующий краевой задаче, и написать условия
его экстремума, то относительно величин $C_{ij}$ получим систему
конечно-разностных уравнений. Другие конечно-разностные схемы получаем
аналогично, используя иные финитные функции (например, сплайны Шенберга
второго и третьего порядков) (см. \S\,15 гл.~3), что соответственно
приводит к менее разреженным матрицам систем алгебраических уравнений.
В этом смысле ничем по существу не отличается от сеточных методов и метод
конечных элементов, в котором функции $\varphi_k(x)$ также локальные сплайны
(в большинстве случаев дефектные, т.\,е. порядок используемых в них
полиномов превосходит их гладкость больше чем на единицу), однако их
носители могут иметь более сложную геометрическую форму, разную в различных
подобластях рассматриваемой области. Именно вследствие этого
и осуществляется более гибкая привязка к геометрии области.

В классических вариационных методах функции $\varphi_k(x)$ "--- это так
называемые координатные функции, вообще, не финитные, удовлетворяющие, если
это необходимо, краевым условиям. Применение таких функций приводит
к системе алгебраических уравнений, как правило, с полностью заполненными
матрицами. Такие системы обычно плохо обусловлены даже при сравнительно
небольших $N$ (порядка $10^2$) и требуют для хранения элементов матрицы
в памяти ЭВМ $N^2$ ячеек. Кроме того, вычисление коэффициентов в методах
типа Ритца "--- Галеркина или наименьших квадратов требует вычисления
кратных интегралов по всей области, в которой решается задача, и как
следствие "--- больших затрат машинного времени. Таким образом, если взять
за основу один и тот же вариационный принцип (например, минимизировать
функционал Ритца), то метод сеток, конечных элементов и классический
вариационный будут отличаться лишь выбором функций $\varphi_k(x)$. Будут
различными при одном и том же $N$ вычислительные затраты и накопление
погрешностей, что, конечно, очень важно, однако тип сходимости, вообще,
будет один и тот же.

Возможен и другой вариант: функции $\varphi_k(x)$ выбираются одни и те же,
но по-разному определяются постоянные $C_k$. Например, используя одну и ту
же формулу~(1), можно находить постоянные по методу Ритца, а можно, добавив
к функционалу Ритца соответствующие слагаемые, применить метод Куранта
и получить благодаря этому более высокий тип сходимости, чем сходимость по
энергии.

Таким образом, все прямые методы объединяет между собой формула~(1),
а отличие между ними состоит лишь в выборе функций $\varphi_k(x)$
и способах отыскания постоянных $C_k$.
\bookpage{11}%

Каким же условиям должны удовлетворять функции $\varphi_k(x)$, если
проблему рассматривать с позиций теории приближений, вычислительной
математики и системного программирования? В большинстве случаев требуется,
а в остальных желательно, чтобы функции $\varphi_k(x)$ удовлетворяли краевым
условиям рассматриваемой задачи. Основная трудность в решении этого вопроса
связана с тем, что в формулы, представляющие функции $\varphi_k(x)$, должна
некоторым образом включаться информация геометрического характера о форме
границы рассматриваемой области (и форме ее участков для краевых условий
смешанного типа). Как уже отмечалось, в некоторых методах (сеточных, МКЭ)
эта информация учитывается приближенно вследствие локального приближения
узлов сетки к границе, а в некоторых (Рейснера, штрафов) "--- в интегральном
смысле. Это, конечно, накладывает определенный отпечаток на тип сходимости
и величину погрешностей. Тем не менее с этими недостатками можно было бы
мириться при условии, что погрешность в результате приближенного учета
геометрических форм соразмерна с погрешностью самого приближенного метода.
Однако есть еще один весьма важный фактор, который необходимо учитывать:
и в сеточных методах, и в МКЭ составление уравнений для пограничных узлов
трудно формализуемое, что затрудняет решение вопросов автоматизации
программирования.

Рассмотрим, какая ситуация складывается в методе штрафов. В простейшем
варианте этого метода краевая задача (3.263), (3.264) (см. \S\,9 гл.~3)
заменяется задачей о минимуме функционала
\begin{equation}\tag{3}
I=\int_\Omega(Au-f)^2\,d\Omega+\sum_{i=1}^{m}M_i\int_{\partial\Omega_i}(L_iu-\varphi_i)^2\,d\Gamma_i,
\end{equation}
где $M_i>0$ "--- некоторые постоянные. Нетрудно заметить, что минимум $I$
достигается на решении (обобщенном) этой краевой задачи. Однако при
численной реализации метода штрафов не ясно, как следует выбирать
постоянные $M_i$. Ответ можно получить на основе численных экспериментов,
но организация этих экспериментов "--- дело сложное, требующее больших
затрат времени. (Видимо, этим и объясняется то обстоятельство, что
к настоящему времени опубликованы лишь немногие решения реальных краевых
задач, полученные методом штрафов, и касаются они лишь простых случаев.)
Поэтому, если бы появилась возможность строить в конструктивном отношении
достаточно простые функции $\varphi_k(x)$ так, чтобы краевые условия
удовлетворялись точно и в то же время были сохранены все другие
положительные факторы, характерные для упомянутых методов, то это имело бы
важное значение для их дальнейшего развития.

В течение длительного времени задача построения в явном виде функций
$\varphi_k(x)$, точно удовлетворяющих произвольным краевым условиям для
областей сложной формы, казалась практически неразрешимой. Дж.~Ортега
и В.~Рейнболдт [95, с.~36] по этому поводу,
\bookpage{12}%
например, пишут: <<Хотя формально метод Ритца распространяется на
двумерные вариационные задачи и задачи более высокой размерности, здесь
имеется серьезная практическая трудность, состоящая в построении подходящих
базисных функций для общих областей>>.

В гл.~3, 4 показано, что эта проблема может быть решена с помощью
конструктивного аппарата теории $R$-функций. При этом речь идет не просто об
однородной задаче Дирихле, а о краевых условиях самых различных типов для
областей практически произвольной формы. Для линейных краевых условий
получаются формулы вида~(1), а для нелинейных "--- формулы общего вида:
\begin{equation}\tag{4}
u_N=f(x,C_1,\dots,C_N).
\end{equation}

По теории $R$-функций "--- довольно много работ (количество их с каждым
годом возрастает), в том числе монографии [111, 127, 128, 134, 135].
Заметим, что в этой теории речь идет не только об удовлетворении краевым
условиям. С точки зрения теории приближений одного лишь удовлетворения
краевым условиям не достаточно для того, чтобы можно было получать
приближения к точному решению. Задачу можно считать успешно решенной лишь
при обеспечении определенных аппроксимационных свойств функций
$\varphi_k(x)$: их линейные комбинации~(1) должны достаточно хорошо
приближать функции того класса, которому принадлежит точное решение задачи,
т.\,е. функции~(1) должны удовлетворять условиям полноты, а по возможности
и условию аппроксимационной универсальности~[126].

Основное внимание в данной монографии обращено на конструктивную сторону
теории $R$-функций. Именно эта сторона сыграла главную роль при создании
генераторов программ серии <<Поле>> (\S\S\,1, 2 гл.~4), позволивших во много
раз сократить затраты на составление и отладку программ при решении краевых
задач. (Более подробная информация содержится в публикациях Г.\,П.~Манько
[67--71].)

Развитие конструктивного начала в математике всегда происходило наряду
с развитием самой теории, причем вызывалось оно в первую очередь запросами
практики. Например, если бы в свое время не был введен символ $\tg x$,
а каждый раз вместо него писали бы $\dfrac{\sin x}{\cos x}$, то от этого не
пострадала бы ни одна из теорем математики. Но нетрудно представить,
насколько усложнилось бы решение многих прикладных проблем.

Функция $\tg x$ выдержала испытание <<естественным отбором>>. В связи
с появлением ЭВМ темпы этого <<отбора>> неизмеримо возросли. Одна из задач
данной работы "--- показать, что есть новые функции (элементарные
и неэлементарные), которые по широте применения на практике и важности для
развития теории, несомненно, выдержат испытание временем.
